% Generated draft: requires mathematical and source review.
\documentclass[11pt]{article}
\usepackage[margin=25mm]{geometry}
\usepackage{fontspec}
\setmainfont{Latin Modern Roman}
\usepackage{amsmath,amssymb,mathtools,unicode-math}
\setmathfont{Latin Modern Math}
\usepackage{enumitem,needspace,xurl,hyperref}
\setlength{\parindent}{0pt}
\setlength{\parskip}{5pt}
\setlength{\emergencystretch}{3em}
\newcommand{\N}{\mathbb{N}}
\newcommand{\Z}{\mathbb{Z}}
\newcommand{\Q}{\mathbb{Q}}
\newcommand{\R}{\mathbb{R}}
\newcommand{\C}{\mathbb{C}}
\newcommand{\F}{\mathbb{F}}
\newcommand{\A}{\mathbb{A}}
\newcommand{\PP}{\mathbb{P}}
\newcommand{\E}{\mathbb{E}}
\newcommand{\eps}{\varepsilon}
\newcommand{\dd}{\,\mathrm{d}}
\begin{document}
\section*{UnsolvedMath ID: 30006093}
\textbf{Source problem code:} catalogue entry
\subsection{Definitions and mathematical statement}
\textbf{Formulation unavailable.} The supplied record for permanent UnsolvedMath ID \(30006093\), source code \texttt{catalogue entry}, consists of a workshop table-of-contents fragment. It begins with the truncated word ``conjecture'' and a page reference, followed by talk titles, names, and page references. It supplies no determinate mathematical assertion, definitions, hypotheses, domains, boundary conditions, or quantifiers. Consequently, no precise mathematical problem can be formulated from this evidence alone. In particular, the report title does not identify which conjecture or mathematical assertion this record was intended to represent.

The supplied statement status is \texttt{unrecoverable}, and no clean statement is supplied. The registry also labels the record \texttt{open}; that label is retained as supplied metadata, not as a verified assertion that an identifiable mathematical problem remains open. The supplied literature assessment explicitly says that the mathematical status is unresolved because the extraction contains no determinate claim. Recovering a mathematical statement requires identification of the intended passage and its complete formulation; no replacement statement is asserted here.
\subsection{Short English statement}
UnsolvedMath ID \(30006093\) preserves a workshop contents-listing fragment rather than a mathematical problem. Its intended assertion and assumptions cannot be recovered from the supplied evidence. The supplied \texttt{open} label does not establish the status of any particular conjecture.
\subsection{Sources}
\begin{itemize}
\item Local UnsolvedMath registry: supplied complete source record, permanent ID \(30006093\), exact source code \texttt{catalogue entry}, title ``Workshop Contents Listing for Algebraic K-Theory,'' in the collection ``Oberwolfach Reports Open Problems.'' No problem proposer is supplied; the proposed-year field is \(2024\).
\item Supplied report citation: \emph{Arbeitsgemeinschaft: Algebraic \(K\)-Theory and the Telescope Conjecture} (2025), Oberwolfach Reports, DOI 10.4171/owr/2024/47, pp.\ 2751--2804. Supplied source link: \url{https://doi.org/10.4171/owr/2024/47}. The background describes Oberwolfach Reports as documenting mathematical workshops and their open-problem discussions; this general description does not establish that the extracted fragment states a problem.
\item Supplied literature source: ``Oberwolfach Report 47,'' \url{https://doi.org/10.4171/OWR/2024/47}. The literature-check date is exactly \texttt{2026-08-22}. This is supplied dated curation, not independently verified research or source inspection. Its status field says ``open,'' while its assessment says that the contents fragment states no determinate mathematical claim and its status is therefore unresolved. The background requests independent verification of time-sensitive status and cited claims.
\item The supplied statement-verification field reports that the official report confirms the extraction is a contents-listing fragment, that no separate clean statement is promoted, and that the original record remains publicly visible and retained for provenance. This verification claim is reported from the supplied record, not independently confirmed here.
\item Contents-listing attribution, as supplied: Shai Keidar, ``Ambidexterity and Chromatic Cyclotomic Extensions'' (2759); Qingyuan Bai, ``Cyclotomic and polygonic spectra'' (2761); Turner McLaurin and Kirsten Wickelgren, ``K-Theory, Land-Tamme, and Levy'' (2762); the name rendered ``Piotr Pstra\k{ } gowski,'' ``The cyclotomic t-structure'' (2768); Jeremy Hahn, ``Disassembling the disproof'' (2768); Achim Krause, ``Purity and Galois descent'' (2769); Leor Neuhauser, ``Cyclotomic Redshift'' (2774); Gabriel Angelini-Knoll, ``Basic Examples of Lichtenbaum--Quillen'' (2775); Lennart Meier, ``Adams operations on BP hni'' (2779); John Rognes, ``Lichtenbaum--Quillen for truncated Brown--Peterson spectra'' (2781); Neil P. Strickland, ``THH of cochains on the circle'' (2783); Liam Keenan, ``Locally unipotent Z-actions'' (2786); Christoph Winges, ``Almost compactness of cyclotomic spectra'' (2791); Maxime Ramzi, ``Asymptotic constancy I: The Dehn twist'' (2794); Sil Linskens, ``Asymptotic Constancy II'' (2796). These are contents entries, not established attributions of a conjecture. The initial truncated ``conjecture'' entry refers to page 2755 without a recoverable title or attribution.
\end{itemize}
\end{document}
